% ============================================================
%  Geometria Espacial — Apostila Premium | PratiqueJá
%  Engine: XeLaTeX
% ============================================================
\documentclass[12pt]{article}
\usepackage{../../pratiqueja}
\usepackage{tikz}

\newcommand{\hlm}[1]{{\color{darkblue}#1}}

\setsubject{Geometria Espacial}

\begin{document}
\pagenumbering{arabic}
\setstretch{1.2}
\color{bodytext}

\heroheader{Geometria Espacial}%
           {Prismas, pirâmides, cilindro, cone e esfera}%
           {Matemática · Avançado}

\setlength{\columnsep}{18pt}
\setlength{\columnseprule}{0.5pt}
\def\columnseprulecolor{\color{exborder}}

\begin{multicols}{2}

\TW{Def}{darkblue}{Prismas}{Área total e volume}

\regra{
Um \hl{prisma} é um sólido com duas bases congruentes e paralelas (polígonos) ligadas por faces laterais retangulares. O nome vem do formato da base: triangular, quadrangular, hexagonal etc.
}

Sendo $P_{\text{base}}$ o perímetro da base e $h$ a altura:

\formula{$A_{\text{lat}} = P_{\text{base}} \times h$}
\formula{$A_{\text{total}} = A_{\text{lat}} + 2 A_{\text{base}}$}
\formula{$V = A_{\text{base}} \times h$}

\exemp{
\textbf{Caixa retangular $5 \times 3 \times 4\,\text{cm}$:}

$A_{\text{base}} = 15\,\text{cm}^2$ \quad $P_{\text{base}} = 16\,\text{cm}$

$A_{\text{lat}} = 16 \times 4 = 64\,\text{cm}^2$

$A_{\text{total}} = 64 + 30 = 94\,\text{cm}^2$

$\hlm{V = 15 \times 4 = 60\,\text{cm}^3}$

\textbf{Prisma triangular equilátero, $a{=}6\,\text{cm}$, $h{=}10\,\text{cm}$:}

$A_{\text{base}} = \dfrac{\sqrt{3}}{4} \times 36 = 9\sqrt{3}\,\text{cm}^2$

$P_{\text{base}} = 18\,\text{cm}$ \quad $A_{\text{lat}} = 180\,\text{cm}^2$

$\hlm{V = 9\sqrt{3} \times 10 = 90\sqrt{3} \approx 155{,}9\,\text{cm}^3}$
}

\TW{Def}{darkblue}{Pirâmides}{Área lateral, área total e volume}

\regra{
Uma \hl{pirâmide} tem uma base poligonal e faces laterais triangulares que convergem para o \hl{ápice}. A \hl{apótema} $a_p$ é a altura de uma face lateral triangular.
}

\formula{$A_{\text{lat}} = \dfrac{P_{\text{base}} \times a_p}{2}$}
\formula{$A_{\text{total}} = A_{\text{lat}} + A_{\text{base}}$}
\formula{$V = \dfrac{A_{\text{base}} \times h}{3}$}

\nota{O volume da pirâmide é sempre $\mathbf{\frac{1}{3}}$ do volume do prisma de mesma base e altura — muito cobrado no ENEM.}

\exemp{
\textbf{Pirâmide quadrada, lado $4\,\text{cm}$, $h{=}3\,\text{cm}$:}

Apótema da base (centro ao meio do lado): $2\,\text{cm}$

$a_p = \sqrt{3^2 + 2^2} = \sqrt{13}\,\text{cm}$

$A_{\text{lat}} = \dfrac{16 \times \sqrt{13}}{2} = 8\sqrt{13} \approx 28{,}8\,\text{cm}^2$

$\hlm{V = \dfrac{16 \times 3}{3} = 16\,\text{cm}^3}$
}

\columnbreak

\TW{Cil}{darkblue}{Cilindro}{Área lateral, área total e volume}

\regra{
O \hl{cilindro reto circular} tem duas bases circulares de raio $r$ e altura $h$. A área lateral corresponde ao retângulo obtido ao "desenrolar" a superfície (base $2\pi r$, altura $h$).
}

\begin{center}\vspace{-4pt}
\begin{tikzpicture}[scale=0.65]
  % cilindro
  \draw[darkblue,thick] (0,0) ellipse (1 and 0.3);
  \draw[darkblue,thick] (-1,0) -- (-1,2);
  \draw[darkblue,thick] (1,0) -- (1,2);
  \draw[darkblue,thick] (0,2) ellipse (1 and 0.3);
  \draw[black!70,<->] (1.15,0) -- node[right,font=\footnotesize]{$h$} (1.15,2);
  \draw[black!70,<->] (0,2.45) -- node[above,font=\footnotesize]{$r$} (1,2.45);
  % cone
  \draw[darkblue,thick] (3,0) ellipse (1 and 0.3);
  \draw[darkblue,thick] (2,0) -- (3,2);
  \draw[darkblue,thick] (4,0) -- (3,2);
  \draw[black!70,<->] (4.15,0) -- node[right,font=\footnotesize]{$h$} (4.15,2);
  \draw[black!70,dashed] (3,0) -- node[below,font=\footnotesize]{$r$} (4,0);
  \node[font=\footnotesize,darkblue] at (0,-0.7) {Cilindro};
  \node[font=\footnotesize,darkblue] at (3,-0.7) {Cone};
\end{tikzpicture}
\vspace{-4pt}\end{center}

\formula{$A_{\text{lat}} = 2\pi r h$}
\formula{$A_{\text{total}} = 2\pi r\,(h + r)$}
\formula{$V = \pi r^2 h$}

\exemp{
\textbf{$r = 3\,\text{cm}$, $h = 10\,\text{cm}$:}

$A_{\text{lat}} = 2\pi \cdot 3 \cdot 10 = 60\pi\,\text{cm}^2$

$A_{\text{total}} = 2\pi \cdot 3 \cdot 13 = 78\pi\,\text{cm}^2$

$\hlm{V = \pi \cdot 9 \cdot 10 = 90\pi\,\text{cm}^3 \approx 282{,}7\,\text{cm}^3}$

\textbf{$V = 200\pi\,\text{cm}^3$, $r = 5\,\text{cm}$. Qual $h$?}

$200\pi = \pi \cdot 25 \cdot h \Rightarrow \hlm{h = 8\,\text{cm}}$
}

\TW{Con}{darkblue}{Cone}{Geratriz, área lateral e volume}

\regra{
O \hl{cone reto circular} tem base circular de raio $r$, altura $h$ e \hl{geratriz} $g$ (distância do ápice à circunferência da base). Pelo Teorema de Pitágoras:
\[
  g = \sqrt{r^2 + h^2}
\]
}

\formula{$A_{\text{lat}} = \pi r g$}
\formula{$A_{\text{total}} = \pi r\,(g + r)$}
\formula{$V = \dfrac{\pi r^2 h}{3}$}

\nota{Assim como a pirâmide, $V_{\text{cone}} = \dfrac{1}{3}\,V_{\text{cilindro}}$ de mesma base e altura.}

\exemp{
\textbf{$r = 6\,\text{cm}$, $h = 8\,\text{cm}$:}

$g = \sqrt{36 + 64} = \sqrt{100} = 10\,\text{cm}$

$A_{\text{lat}} = \pi \cdot 6 \cdot 10 = 60\pi\,\text{cm}^2$

$A_{\text{total}} = \pi \cdot 6 \cdot 16 = 96\pi\,\text{cm}^2 \approx 301{,}6\,\text{cm}^2$

$\hlm{V = \dfrac{\pi \cdot 36 \cdot 8}{3} = 96\pi\,\text{cm}^3 \approx 301{,}6\,\text{cm}^3}$
}

\TW{Esf}{badgepurple}{Esfera}{Área da superfície e volume}

\regra{
A \hl{esfera} é o conjunto de todos os pontos do espaço à distância $r$ (raio) de um ponto fixo (centro). Toda secção plana pelo centro é um \hl{círculo máximo}.
}

\begin{center}\vspace{-4pt}
\begin{tikzpicture}[scale=0.7]
  \draw[darkblue,thick] (0,0) circle (1.2);
  \draw[darkblue,dashed] (0,0) ellipse (1.2 and 0.35);
  \draw[exborder,<->] (0,0) -- node[above,font=\footnotesize]{$r$} (0.85,0.85);
\end{tikzpicture}
\vspace{-4pt}\end{center}

\formula{$A_{\text{sup}} = 4\pi r^2$}
\formula{$V = \dfrac{4\pi r^3}{3}$}

\exemp{
\textbf{$r = 3\,\text{cm}$:}

$A_{\text{sup}} = 4\pi \cdot 9 = 36\pi\,\text{cm}^2 \approx 113{,}1\,\text{cm}^2$

$\hlm{V = \dfrac{4\pi \cdot 27}{3} = 36\pi\,\text{cm}^3 \approx 113{,}1\,\text{cm}^3}$

\textbf{$V = 288\pi\,\text{cm}^3$. Qual o raio?}

$\dfrac{4\pi r^3}{3} = 288\pi \Rightarrow r^3 = 216$

$\hlm{r = \sqrt[3]{216} = 6\,\text{cm}}$
}

\nota{\textbf{Relação:} $A_{\text{sup}} = 4\pi r^2$ é exatamente 4 vezes a área do círculo máximo ($\pi r^2$).}

\TW{Rev}{badgepurple}{Sólidos de Revolução}{Cilindro e cone gerados por rotação}

\regra{
Um \hl{sólido de revolução} é gerado pela rotação de uma figura plana em torno de um eixo:
\begin{itemize}
  \item Retângulo girado em torno de um lado → \textbf{cilindro}
  \item Triângulo retângulo girado em torno de um cateto → \textbf{cone}
\end{itemize}
}

\exemp{
\textbf{Retângulo de lados $r$ e $h$ girado em torno de $h$:}

O lado $r$ descreve um círculo de raio $r$.

$\hlm{V = \pi r^2 h}$ \quad $\hlm{A_{\text{lat}} = 2\pi r h}$

\textbf{Triângulo retângulo, catetos $r$ e $h$, girado em torno de $h$:}

O cateto $r$ gera a base circular; a hipotenusa gera a superfície cônica.

$\hlm{V = \dfrac{\pi r^2 h}{3}}$ \quad $\hlm{g = \sqrt{r^2 + h^2}}$

\textbf{Retângulo $3 \times 5\,\text{cm}$ girado em torno do lado $5\,\text{cm}$:}

Cilindro: $r = 3\,\text{cm}$, $h = 5\,\text{cm}$

$V = \pi \cdot 9 \cdot 5 = 45\pi\,\text{cm}^3$

$\hlm{A_{\text{total}} = 2\pi \cdot 3 \cdot 8 = 48\pi\,\text{cm}^2}$
}

\end{multicols}

\end{document}
